% =====================================================================
%  CAPÍTULO NUEVO: RESULTADOS Y DISCUSIÓN
%  Insertar en main.tex ENTRE 04-desarrollo y 05-conclusiones (que pasa a
%  ser el capítulo 6). Reutiliza las tablas ya existentes de produccion.tex
%  (tab:metricas-reconocimiento, tab:umbral-reconocimiento,
%  tab:metricas-rutinas); no duplica cifras.
%  Todas las cifras provienen de ml/results/*.json o del texto vigente de
%  la memoria. Lo no medido se marca con \pendiente{}.
% =====================================================================
\chapter{RESULTADOS Y DISCUSIÓN}
\label{cap:resultados}

Este capítulo reúne los resultados del proyecto, los contrasta con los
objetivos específicos y con los antecedentes citados, y declara las amenazas a
la validez que limitan su interpretación. Las cifras de inteligencia artificial
proceden de los archivos de resultados del repositorio
(\archivo{ml/results/}); cuando una medida no se realizó se indica de forma
explícita.

\section{Resultados por objetivo específico}
\label{sec:res-objetivos}

La Tabla~\ref{tab:res-objetivos} resume, para cada objetivo específico, lo que
se construyó, la evidencia disponible y el grado de cumplimiento.

\begin{table}[H]
\centering
\caption{Cumplimiento de los objetivos específicos}
\label{tab:res-objetivos}
\small
\begin{tabularx}{\textwidth}{|L{2.9cm}|Y|L{3.5cm}|}
\hline
\rowcolor{encabezado} \textbf{Objetivo} & \textbf{Resultado} & \textbf{Cumplimiento} \\
\hline
OE1. Guiar sobre el uso correcto de máquinas &
Catálogo de 35 tipos de máquinas con guía de uso, errores comunes y
prevención de lesiones; escáner que superpone la ficha sobre la cámara y
reconocedor por imagen con acierto top-1 de 38,5\,\% en campo simulado
(Tabla~\ref{tab:metricas-reconocimiento}). & Parcial: guía completa;
reconocimiento modesto y sin prueba en gimnasios reales. \\ \hline
OE2. Instruir en la ejecución correcta &
95 animaciones de ejecución con músculo objetivo y dificultad, asociadas a las
máquinas; entrenamiento guiado paso a paso. &
Parcial: contenido instructivo, sin verificación automática de la postura. \\ \hline
OE3. Rutinas personalizadas y estructuradas &
Generador basado en reglas con recomendador de TensorFlow Lite y respaldo por
reglas (Tabla~\ref{tab:metricas-rutinas}). &
Cumplido; la calidad de las rutinas para el usuario no se evaluó. \\ \hline
OE4. Herramientas de planificación &
Plan semanal por días disponibles y registro de la racha de entrenamiento. &
Cumplido. \\ \hline
OE5. Registro de ingesta calórica &
Búsqueda de alimentos en FoodData Central mediante una Cloud Function con
cuotas por usuario, cálculo de porciones y del IMC. &
Cumplido dentro del alcance (conteo calórico, no consejo nutricional). \\ \hline
OE6. Seguimiento del progreso físico &
Registro de peso, estatura y edad con cálculo y clasificación del IMC,
sincronizado con Firestore. & Cumplido. \\ \hline
OE7. Reportes personalizados &
Reportes de progreso y metas. &
Parcial: \pendiente{ejecutar y registrar la prueba de generación del PDF en un
dispositivo físico (Anexo~C).} \\
\hline
\end{tabularx}
\fuente{Elaboración propia a partir del capítulo~\ref{cap:desarrollo} y del
repositorio del proyecto.}
\end{table}

\section{Resultados de los modelos de inteligencia artificial}
\label{sec:res-ia}

\subsection{Reconocimiento de máquinas}

Con 399 fotografías de 35 clases, validación cruzada de cinco pliegues y 30
aumentos por fotografía, el modelo acierta la máquina en el 38,5\,\% de los
casos en el escenario de campo simulado y en el 61,4\,\% entre sus tres
primeras sugerencias, frente al 2,9\,\% que lograría el azar
(Tabla~\ref{tab:metricas-reconocimiento}). Al fijar el umbral de confianza en
0,8, la aplicación acepta automáticamente solo el 15,8\,\% de las imágenes, con
una precisión del 85,4\,\% (Tabla~\ref{tab:umbral-reconocimiento}).

El análisis por clase muestra que el rendimiento es muy desigual: de las 33
clases evaluables, 25 quedan por debajo del 50\,\% de acierto en el campo
simulado, mientras que la caminadora alcanza el 92,9\,\% y la bicicleta
estática el 66,7\,\%. Las confusiones se concentran entre máquinas de aspecto
similar: el 50\,\% de las muestras de validación de la máquina combinada de
extensión y curl de piernas se clasificó como curl femoral, y el 26,5\,\% de
las de press de hombro como press de pecho
(\archivo{matriz\_confusion\_campo} en
\archivo{recognition\_metrics.json}). Este patrón es coherente con un conjunto
de entrenamiento pequeño y con imágenes de catálogo, y sugiere que la mejora
proviene de más fotografías tomadas en gimnasios reales antes que de una
arquitectura distinta (Sandler et~al., 2018).

\subsection{Recomendador de rutinas}

El recomendador reproduce las reglas con una exactitud del 99,8\,\% en datos
sintéticos y del 97,1\,\% en una cuadrícula exhaustiva de 28\,512 perfiles
(Tabla~\ref{tab:metricas-rutinas}). Dado que las etiquetas se generaron con las
mismas reglas, estas cifras miden la fidelidad del modelo y no demuestran que
sus rutinas sean mejores que las del generador por reglas, cuya base de
conocimiento se apoya en la literatura de prescripción del entrenamiento de
fuerza (Kraemer \& Ratamess, 2004). La ventaja práctica del modelo es
arquitectónica: demuestra la integración de inferencia local en el
dispositivo con respaldo por reglas.

\section{Discusión}
\label{sec:discusion}

\subsection{Frente al problema de adherencia}

La introducción parte del abandono temprano en los gimnasios. La evidencia
revisada indica que menos del 40\,\% de los socios mantiene un ejercicio regular
a largo plazo (Gjestvang et~al., 2021), por lo que las barreras de orientación,
estructura y seguimiento que abordan los objetivos específicos son pertinentes.
Sin embargo, este proyecto no midió si la aplicación mejora la adherencia: la
hipótesis queda sin contrastar y requiere un estudio con usuarios reales,
consentimiento informado y un periodo de seguimiento (sección~\ref{sec:recomendaciones}).

\subsection{Frente a la visión artificial para la ejecución de ejercicios}

Silva Amado (2024) propone un sistema de visión artificial que estima la pose
con MediaPipe, calcula ángulos articulares y los compara con umbrales para
identificar la ejecución de ejercicios, y lo evalúa de forma cualitativa. Los
enfoques son complementarios: aquel verifica \emph{cómo} se ejecuta un
movimiento; el presente trabajo reconoce \emph{con qué máquina} se entrena,
genera la rutina y registra el progreso, y aporta métricas cuantitativas
reproducibles. En consecuencia, la instrucción de la ejecución correcta (OE2)
se resuelve aquí con contenido guiado y no con verificación automática de la
postura, que se identifica como línea de trabajo futuro.

\subsection{Sobre la realidad aumentada}

El escáner superpone información digital sobre la imagen en vivo de la cámara,
lo que cumple la primera característica de la realidad aumentada descrita por
Azuma (1997) (combinar lo real y lo virtual) y de forma parcial la segunda
(interacción en tiempo real), pero no la tercera (registro en tres
dimensiones), porque no utiliza ARCore ni seguimiento de superficies. Por ello,
en esta memoria el término se usa en el sentido de \emph{superposición
informativa sobre la cámara}, y el anclaje tridimensional se plantea como
trabajo futuro (sección~\ref{sec:recomendaciones}).

\section{Amenazas a la validez y limitaciones}
\label{sec:amenazas}

\begin{itemize}
    \item \textbf{Validez de constructo del reconocedor.} Las 399 fotografías
    provienen de catálogos y no de gimnasios; el escenario de campo simulado
    aproxima, pero no reemplaza, las condiciones reales de luz, encuadre y
    oclusión.
    \item \textbf{Circularidad del recomendador.} Los datos sintéticos están
    etiquetados por las reglas que el modelo imita; la métrica no evalúa la
    utilidad para el usuario.
    \item \textbf{Ausencia de evaluación con usuarios.} No se midieron
    usabilidad, satisfacción ni adherencia. \pendiente{ejecutar la prueba de
    usabilidad con usuarios y registrar los resultados.}
    \item \textbf{Cobertura de pruebas.} Las pruebas automatizadas cubren la
    lógica pura y las reglas de seguridad; no cubren la interfaz ni se midieron
    la latencia y la exactitud de los modelos en un dispositivo físico.
    \item \textbf{Plataforma y conectividad.} El sistema funciona en Android y
    requiere conexión a internet; la persistencia sin conexión no está
    configurada.
    \item \textbf{Alcance nutricional.} El registro de alimentos se limita al
    conteo calórico y no sustituye la supervisión de un nutricionista.
\end{itemize}
